\documentclass[sigconf,review,anonymous]{acmart}

\setcopyright{none}
\settopmatter{printfolios=true}
\renewcommand\footnotetextcopyrightpermission[1]{}
\acmConference[AgenticDev '26]{International Workshop on Agentic AI for Next-Generation Software Development}{October 12, 2026}{Munich, Germany}
\acmBooktitle{International Workshop on Agentic AI for Next-Generation Software Development (AgenticDev '26), October 12, 2026, Munich, Germany}
\acmDOI{}
\acmISBN{}

\usepackage{booktabs}
\usepackage{balance}
\usepackage{xspace}

\newcommand{\vh}{\textsc{VeriHarness}\xspace}
\newcommand{\rawdiag}{\textsc{RawDiag}\xspace}
\newcommand{\samenl}{\textsc{SameNL}\xspace}
\newcommand{\locobs}{\textsc{LocObs}\xspace}
\newcommand{\typedfields}{\textsc{TypedFields}\xspace}

\begin{document}
\setlength{\emergencystretch}{1em}

\title{VeriHarness: What Should a Validator Tell an Agent?}

\author{Anonymous Author(s)}
\affiliation{\institution{Anonymous Institution}}

\begin{abstract}
We study iterative software systems in which a \emph{model} (an LLM) generates a
candidate and a \emph{validator} (external checking code) returns pass or a
failure. We call the repeated model--validator--feedback sequence an
\emph{agent loop}. What should the validator return to the next call? \vh is a
code-as-harness prototype in which LLMs generate every leaf action while code
controls validation, acceptance, budgets, and traces.
We compare raw diagnostics with feedback that exposes a failure location,
expected alternatives, and the observed value.

Our primary experiment contains 400 oracle-blind rows: four repair policies,
50 paired TextWorld games, two model families, and a four-call cap. Structured
verifier feedback raises success from 14/50 to 36/50 for
Qwen2.5-Coder-14B (+44 points,
95\% paired-bootstrap interval 28--60) and from 8/50 to 29/50 for
Llama-3.1-8B (+42 points, 28--56). However, rendering the same information as
untyped prose yields 35/50 and 29/50. Typed structure itself changes success by
only +2 and 0 points. Removing expected alternatives largely removes the gain.
The ordering persists across call budgets and sampled decoding, but not on a
small public-test HumanEval scope check. Across 880 rows and 2,652 LLM calls,
the evidence supports a narrow principle: expose explicit repair fields; do not
attribute the benefit to JSON typing alone.
\end{abstract}

\begin{CCSXML}
<ccs2012>
 <concept>
  <concept_id>10011007.10011006.10011060.10011061</concept_id>
  <concept_desc>Software and its engineering~Software testing and debugging</concept_desc>
  <concept_significance>500</concept_significance>
 </concept>
 <concept>
  <concept_id>10010147.10010178.10010219.10010221</concept_id>
  <concept_desc>Computing methodologies~Intelligent agents</concept_desc>
  <concept_significance>300</concept_significance>
 </concept>
</ccs2012>
\end{CCSXML}

\ccsdesc[500]{Software and its engineering~Software testing and debugging}
\ccsdesc[300]{Computing methodologies~Intelligent agents}
\keywords{LLM agents, verifier feedback, repair, agent harnesses, TextWorld}

\maketitle

\section{Introduction}

We use three terms throughout. A \emph{model} is the LLM called to generate one
candidate action or artifact. A \emph{validator} is external code that checks a
candidate and returns either pass or a failure description. An \emph{agent loop}
is the bounded sequence model $\rightarrow$ validator $\rightarrow$ feedback,
repeated until validation passes or the call budget is exhausted. In software
development, such a loop may write a patch or plan, invoke a test or execution
environment, and react to failure. ReAct, Reflexion, Self-Refine, and
self-debugging establish the value of iterative feedback
\cite{yao2023react,shinn2023reflexion,madaan2023selfrefine,chen2024selfdebug}.
Coding agents also benefit from purpose-built interfaces
\cite{yang2024sweagent}. Yet one interface remains underspecified: \emph{what
exactly should a validator tell the next model call?}

Suppose command 1 in a plan is invalid. A raw retry sees ``Command 1 is not
admissible.'' A repair interface can additionally expose the observed command
and the commands admissible at that state. This may help because the information
is actionable, because the locus is explicit, or because a typed record is
easier to parse. Prior evidence did not separate these explanations.

We use \vh, a code-as-harness prototype, to ask:
\begin{description}
  \item[RQ1.] Does structured verifier feedback with expected alternatives improve repair
  over raw diagnostics under equal compute?
  \item[RQ2.] Is any gain caused by typed syntax, structured content, or both?
\end{description}

Replications across models, call budgets, and decoding, plus a public-test code
scope check, test the robustness and boundary of RQ1--RQ2; they are not a third
research question.

Our answer is deliberately narrow. Structured verifier feedback produces a large,
replicated gain on verifier-rich TextWorld plans. Typed JSON and
information-matched prose have indistinguishable success. A location without
expected alternatives is insufficient. The main contribution is therefore an
empirical decomposition of validator feedback, plus a reproducible harness for
testing it.

\section{Background and System}

Reflexion stores verbal feedback, while Self-Refine lets an LLM produce and use
its own critique \cite{shinn2023reflexion,madaan2023selfrefine}. Execution-guided
methods reuse failed programs and test messages \cite{chen2024selfdebug}.
SWE-bench and SWE-agent show the importance of executable evaluation and
agent--computer interfaces \cite{jimenez2024swebench,yang2024sweagent}. \vh
studies a complementary boundary: deterministic validator output passed to the
next probabilistic leaf.

\subsection{Code as the Control Plane}

\vh separates generation from acceptance (Fig.~\ref{fig:loop}). A principal
orchestrator builds a compact state pack and dispatches bounded LLM workers.
Each worker returns a structured candidate. External gates check schema,
artifacts, and environment behavior. Only gates accept a candidate; a worker's
self-assessment cannot do so. Every prompt, raw response, parsed output, gate
result, model setting, and elapsed time is persisted before the next call.

Code is preferable here to prose as the \emph{control plane}. A prompt-only
loop asks the model to remember workflow state, count calls, interpret its own
completion, and carry an expanding history in context. Those duties compete
with candidate generation, and model-declared completion cannot independently
validate the model's output. Code instead keeps durable state outside the
prompt, enforces budgets and acceptance gates deterministically, and permits the
same tasks and failures to be replayed across policies. The model receives only
the compact task state and current repair feedback. This is an engineering
rationale for \vh's architecture; the experiment compares feedback policies,
not code control against prompt-only control.

\begin{figure}[t]
\centering
\fbox{\begin{minipage}{0.92\columnwidth}
\centering\small
Task + compact state $\rightarrow$ LLM worker $\rightarrow$ candidate\\[2pt]
$\downarrow$\hspace{98pt}$\downarrow$\\[-2pt]
budget + trace\hspace{24pt}external gates\\[2pt]
$\uparrow$\hspace{101pt}$\downarrow$\\[-2pt]
repair feedback $\leftarrow$ failure encoder
\end{minipage}}
\caption{LLMs generate leaf candidates; code controls state, validation,
acceptance, budgets, and traces.}
\Description{A loop from task and compact state to an LLM worker, candidate,
external gates, failure encoder, and repair feedback.}
\label{fig:loop}
\end{figure}

The failure encoder maps gate-specific details to a common interface. For an
invalid TextWorld action it can expose a stable label, the command index, the
admissible actions in environment order, and the rejected action. Oracle-blind
evaluation excludes hidden post-hoc outcomes from repair feedback.

\subsection{Feedback Policies}

All policies receive the same task state, rejected answer, gate, model, and
call cap. They differ only after a failed candidate (Table~\ref{tab:policies}).

\begin{table}[t]
\caption{Feedback controls used in the primary experiment.}
\label{tab:policies}
\small
\begin{tabular}{@{}p{0.22\columnwidth}p{0.70\columnwidth}@{}}
\toprule
Policy & Feedback after failure \\
\midrule
\rawdiag & Raw validation message. \\
\samenl & Exact location, expected alternatives, and observed value in ordinary
prose, without typed labels or keys. \\
\locobs & Typed label, location, and observed value; no expected alternatives. \\
\typedfields & Typed label, location, expected alternatives, and observed value. \\
\bottomrule
\end{tabular}
\end{table}

\samenl and \typedfields contain the same actionable field values; their
contrast isolates representation. \locobs and \typedfields share typed
structure; their contrast isolates expected alternatives. \rawdiag versus
\samenl measures the value of adding structured semantic content to the raw
message. Here, \emph{structured} means explicit repair components---location,
expected alternatives, and observation---whether rendered as prose or JSON.

\section{Study Design}

\textbf{Primary tasks.} We generate 50 fixed TextWorld games
\cite{cote2018textworld}, with seeds 20261051--20261100. Each has three rooms,
four objects, and a two-step quest. A candidate is a JSON plan of at most four
commands, executed from fresh state. Success requires a terminal win. Invalid
actions return their index and up to the first 12 admissible commands in the
environment's deterministic order; we neither rank nor randomize this list.
The primary matrix is $50$ games $\times$ four policies $\times$ two models,
or 400 rows.

\textbf{Models and decoding.} We use Qwen2.5-Coder-14B-Instruct-AWQ
\cite{hui2024qwen25coder} on an NVIDIA L4 and
Meta-Llama-3.1-8B-Instruct-AWQ-INT4
\cite{grattafiori2024llama3} on a T4, served by vLLM. Primary decoding is greedy with
512 output tokens and a four-call cap. Model conditions are never mixed across
serving backends.

\textbf{Robustness lanes.} First, on the same 15-game subset and both models,
we compare \rawdiag and \typedfields at budgets 2, 4, 6, and 8. Budget-4 rows
are reused from the primary matrix. Second, Qwen runs 20 games at temperature
0.3, top-$p$ 0.9, and three inference seeds; this lane compares \rawdiag,
\samenl, and \typedfields. Third, a 15-task HumanEval scope check exposes one
deterministically selected public assertion to the repair gate while retaining
the complete official tests for post-hoc success. These lanes contribute 480
additional rows, for 880 total and 2,652 realized LLM calls.

\textbf{Analysis.} Games are paired by seed. Primary success differences use
10,000 paired bootstrap samples and two-sided exact McNemar tests. We apply Holm
correction to four planned contrasts within each model. Sampled-decoding
intervals resample 20 game clusters, retaining all three repeats; $p$-values use
100,000 game-level sign flips and Holm correction. Budget and HumanEval lanes
are small robustness checks, so we emphasize counts and intervals rather than
dichotomizing their $p$-values. Saved-transcript whitespace counts are a prompt
size proxy, not tokenizer or billing tokens.

\section{Results}

\subsection{Structured Verifier Feedback Produces the Gain}

Table~\ref{tab:primary} reports terminal success. \typedfields improves over
\rawdiag by 44 points for Qwen (95\% interval 28--60; Holm-adjusted exact
$p=3.15\times10^{-5}$) and 42 points for Llama (28--56;
$p=3.81\times10^{-6}$). The result therefore replicates across a coder model
and a smaller general model.

\begin{table}[t]
\caption{Primary TextWorld results: 50 paired games per model, budget 4.}
\label{tab:primary}
\small
\begin{tabular}{@{}lrrrr@{}}
\toprule
& \multicolumn{2}{c}{Qwen-14B/L4} & \multicolumn{2}{c}{Llama-8B/T4} \\
Policy & Solved & Calls & Solved & Calls \\
\midrule
\rawdiag & 14 (28\%) & 164 & 8 (16\%) & 179 \\
\locobs & 18 (36\%) & 155 & 9 (18\%) & 174 \\
\samenl & 35 (70\%) & 147 & 29 (58\%) & 161 \\
\typedfields & \textbf{36 (72\%)} & \textbf{130} & \textbf{29 (58\%)} & \textbf{149} \\
\bottomrule
\end{tabular}
\end{table}

The ablations identify the active information (Table~\ref{tab:causal}).
\samenl beats \rawdiag by 42 points on both models, with no baseline-only
successes. In contrast, \typedfields versus the information-matched \samenl is
only +2 points for Qwen and 0 for Llama; both intervals include zero and both
adjusted $p$-values are 1. Typed structure does not explain success.

Expected alternatives do. \locobs, which retains typed label, location, and
observation but omits alternatives, is near \rawdiag. Adding alternatives raises
success by 36 points for Qwen and 40 for Llama. Final invalid-action episodes
fall from 35 to 2 (raw to same-information prose) for Qwen and from 41 to 4 for
Llama. A location says where repair is needed; admissible alternatives say what
repair is possible.

\begin{table}[t]
\caption{Paired contrasts for Qwen2.5-Coder-14B-Instruct-AWQ and
Meta-Llama-3.1-8B-Instruct-AWQ-INT4. CIs are paired bootstrap; $p_H$ is
Holm-adjusted exact McNemar within model.}
\label{tab:causal}
\small
\begin{tabular}{@{}llrrr@{}}
\toprule
Model & Treatment vs. baseline & $\Delta$ & 95\% CI & $p_H$ \\
\midrule
Qwen & SameNL vs. raw & +42 & [28,56] & $3.8\!\times\!10^{-6}$ \\
Qwen & Typed vs. SameNL & +2 & [$-$8,12] & 1.0 \\
Qwen & Typed vs. LocObs & +36 & [20,52] & $2.4\!\times\!10^{-4}$ \\
Llama & SameNL vs. raw & +42 & [28,56] & $3.8\!\times\!10^{-6}$ \\
Llama & Typed vs. SameNL & 0 & [$-$12,12] & 1.0 \\
Llama & Typed vs. LocObs & +40 & [26,54] & $2.2\!\times\!10^{-5}$ \\
\bottomrule
\end{tabular}
\end{table}

Typing has a smaller efficiency effect. At equal success, \typedfields uses 17
fewer calls than \samenl for Qwen (paired mean $-0.34$ calls/game, interval
$-0.60$ to $-0.06$) and 12 fewer for Llama ($-0.24$, $-0.44$ to $-0.04$).
Average retry prompts are similar: 701 versus 716 words for Qwen and 706 versus
722 for Llama. Total realized prompt words are 14\% and 10\% lower for typed
feedback because it reaches stopping conditions sooner, not because each repair
prompt is dramatically smaller.

\subsection{Budgets and Sampled Decoding}

The budget result is progressive rather than a universal step at call four
(Table~\ref{tab:budget}). On 15 paired games, typed-minus-raw gains rise from 20
and 27 points at budget 2 to 47 points for both models at budget 4. Qwen then
plateaus near 47--53 points; Llama reaches 60 points at budgets 6 and 8. More
calls do not rescue raw retry: its Qwen count is 4/15 at budgets 4--8 and its
Llama count is 2/15. The repair interface uses extra compute more effectively.

\begin{table}[t]
\caption{Budget sensitivity on the same 15 games. Cells are raw/typed wins.}
\label{tab:budget}
\small
\begin{tabular}{@{}lrrrr@{}}
\toprule
Model & $B=2$ & $B=4$ & $B=6$ & $B=8$ \\
\midrule
Qwen & 5/8 & 4/11 & 4/11 & 4/12 \\
Llama & 1/5 & 2/9 & 2/11 & 2/11 \\
\bottomrule
\end{tabular}
\end{table}

Sampled Qwen decoding gives the same ordering. Across seeds 3101--3103,
\rawdiag solves 6, 6, and 5 of 20 games; \samenl solves 14, 13, and 14; and
\typedfields solves 14, 16, and 14. Clustered across games, typed versus raw is
+45 points (95\% interval 23--67; Holm $p=.0053$), while typed versus
information-matched prose is +5 points ($-3$ to 15; $p=.499$). The primary
conclusion is not an artifact of greedy repetition.

\subsection{When the Validator Cannot Expose the Hidden Failure}

The HumanEval lane tests a boundary of the mechanism, not another positive
benchmark. Its repair validator executes one visible assertion; the complete
official suite remains hidden and is used only for post-hoc success. Qwen's
initial candidate passes the public assertion on all 15 tasks under every
policy, so no repair feedback is ever issued. Fourteen pass the hidden suite.
The remaining hidden failure is therefore invisible to every repair policy,
which necessarily produces the identical 14/15 outcome.

Llama creates a few repair opportunities: 4 initial public failures under raw,
same-information, and location-only feedback, and 5 under typed feedback.
Retrying adds one public-gate success under raw, same-information, and typed
feedback, but none under location-only feedback. Only the repaired raw case also
passes the hidden suite. Final hidden success is 10/15 for raw,
same-information, and location-only feedback and 9/15 for typed feedback.

This result clarifies RQ1's scope. Structured feedback helps only after a
validator observes a repair-relevant failure. It cannot reveal behavior that
the validator never tests, and passing one visible assertion need not predict
hidden correctness. The lane therefore supports a validator-observability
boundary; it does not support a claim of general code-repair improvement.

\section{Implications}

\textbf{Treat validators as repair APIs.} A validator should expose a stable
repair locus, the observed value, and actionable expectations when available.
For software, these may be a file and line, failing assertion, expected type or
behavior, and observed exception. When a test cannot enumerate alternatives,
the interface should report what it knows rather than fabricate them.

\textbf{Content before syntax.} JSON is useful for logging, routing, and
cross-gate tooling, and our typed condition uses modestly fewer calls. But the
same information in prose achieves the same success. Claims that
machine-readable feedback inherently improves reasoning would overstate our
evidence.

\textbf{Budget repair, not repetition.} Increasing a call cap helps only when
later calls receive information that changes the candidate. Raw retry remains
flat after four calls, while structured feedback continues to benefit the
weaker model through six calls.

\section{Boundary Conditions}

These boundaries define the claim rather than form a generic checklist. The
evidence supports structured feedback for deterministic validators that can
expose a repair locus and useful alternatives. It does not establish repository
repair or industrial deployment: the primary tasks are 50 generated games, the
models are two quantized checkpoints on one serving stack, and sampled decoding
uses one model and three seeds. The 15-task HumanEval lane further shows that
feedback policy is ineffective when the visible validator misses the hidden
failure.

The same-information prose control matches location, expected, and observed
values, although wording and delimiters inevitably differ from JSON. Admissible
lists are capped at 12 in deterministic environment order; ranking,
randomization, and longer lists remain untested. Transcript words are a size
proxy, not billed tokens. Real test tools may also be flaky or incomplete,
requiring reruns and quarantine. We use Holm correction for the four primary
contrasts; budget and code results remain small robustness checks. One transient
serving error was quarantined and rerun identically, and only the replacement is
analyzed.

\section{Artifact and Conclusion}

The artifact contains source, deterministic task rules, 880 row-level results,
2,652 per-call traces, model and decoding settings, the quarantined transient
row, and a script that regenerates all tables, 10,000-sample bootstrap intervals,
exact McNemar tests, Holm corrections, and clustered sampling analysis.

Validator feedback is an agent interface, not incidental logging. Across two
models, structured verifier feedback containing location and alternatives more
than doubles TextWorld repair
success over raw diagnostics. Information-matched prose performs equally well,
so typing is not the causal success mechanism. For agentic development, the
practical lesson is precise: expose explicit repair fields first; choose a
typed representation for systems reasons, not because this study shows that
JSON itself makes an LLM reason better.

\balance
\bibliographystyle{ACM-Reference-Format}
\bibliography{references}

\end{document}
